% ============================================================================
\clearpage
\section*{v4.9.0-r21 Source-Only: The Clock Bound Frame by Frame, and a Derivation of \texorpdfstring{\(1/5\)}{1/5}}
\addcontentsline{toc}{section}{v4.9.0-r21 Source-Only: The Clock Bound Frame by Frame, and a Derivation of 1/5}
\label{sec:v490-r21-clock-bound-per-frame}
% ============================================================================

\paragraph{Purpose.}
r20 traced \(\Delta_\gamma<1/5\) to the Hopf point and found that the bound changes away from it.  How it changes depended on
how the chart was carried over.  On 2026-10-22 the author asked to compute the bound separately for each frame, since a
distinction by frame may have to be built in everywhere, and to attempt a concrete derivation.  This layer states five
premises, derives the bound per frame from them, and shows that \(1/5\) follows from one consistency condition.  The script is
\texttt{llm/k60\_delta\_gamma\_je\_frame.py}.  No new prediction is made and nothing is compared with data.  Earlier layers are
not rewritten.

\noindent\textbf{Source-only status.}
This v4.9.0-r21 layer is source-only.  It does not build a PDF, does not assign a new DOI,
and does not authorize a public release.  The v4.9.0 Version DOI
\href{https://doi.org/10.5281/zenodo.23534954}{10.5281/zenodo.23534954} and the Concept DOI
\href{https://doi.org/10.5281/zenodo.19210728}{10.5281/zenodo.19210728} are unchanged by it.

\subsection*{v4.9.0-r21 premises}
\label{subsec:v490-r21-premises}

The window condition is that of r12, \(G(1-\Delta_\gamma)S>16\) with \(S=N+\kappa N(N-1)\).  The doublet state of r17 at
\(x=\bsym^2\) has \(|b|=1/\gamma\), \(\gamma=(1+x)/(1-x)\), \(p_-=x/(1+x)\), \(p_+=1/(1+x)\).
\begin{enumerate}[leftmargin=2em]
\item[(D1)] \emph{Clock gain.}  The photon-frame clock restores the visibility \(|b|\) of the compensator readout to \(1\), so
\(G=1/|b|=|\gamma|\).  This is a readout quantity in the sense of r19 and keeps \(\gL\) (way A).
\item[(D2)] \emph{Pair coupling.}  \(\kappa\) is the content of the partner port.  In the matter frame (\(b_1>0\)) that is \(p_-\).  In
the tachyon frame (\(b_1<0\)) \(J_B\) exchanges the ports (r02, r04), so it is \(p_+\).
\item[(D3)] \emph{Balance.}  In the frame that needs the clock, the partner content of the other packets refills one packet:
\(\kappa(N-1)=1\).  This is the balance condition that singles out \(V3\) in r12.
\item[(D4)] \emph{Packet count.}  \(N\) counts packets and is the same in every frame.
\item[(D5)] \emph{Tachyon description.}  In the tachyon frame the chart is the clock seed of the source,
\((N,\kappa,G)=(5,\tfrac34,1)\), so \(G_T=1\).
\end{enumerate}
(D2) and (D3) follow from r02, r04 and r12.  (D1) gives a mechanism to the gain that r14 calls photon-frame timing.  (D4) and
(D5) are assumptions.  (D5) is taken from the data of the source.

\subsection*{v4.9.0-r21 the derivation}
\label{subsec:v490-r21-derivation}

\paragraph{Theoremlet v4.9.0-r21.1.}
Under (D1)--(D5), with \(N=1+1/p_-=2+1/x\) from (D3), the survival reserves \(1-16/(GS)\) are
\[
\Delta_M^{\max}(x)=\frac{10x^2-5x+1}{(1+x)(1+2x)},\qquad
\Delta_T^{\max}(x)=\frac{1+3x-14x^2}{(1+x)(1+2x)} .
\]
Require that the window's margin be the same whichever frame describes it, \(G_MS_M=G_TS_T\).  This has the single solution
\(x=1/3\), the Hopf point.  There \(N=5\), \(S_M=10\), \(S_T=20\), and
\[
\Delta_\gamma^{\max}=\frac15
\]
in both frames.  Without (D3), with \(N=5\) held fixed, the same condition reads \(4p_-^2-9p_-+2=0\), with the admissible root
\(p_-=1/4\), again the Hopf point.

\paragraph{Proof sketch.}
Substitute (D1)--(D5) and simplify symbolically; solve \(\gamma S_M=S_T\) for \(x>0\).  Checked with a computer algebra system.  \(\square\)

\paragraph{Reading.}
This is a derivation of \(1/5\) from the doublet state and one consistency condition: one physical window must not survive in
one frame and fail in the other.  The integer chart \((5,\tfrac14,2)\) and the seed \((5,\tfrac34,1)\) of the source are then its
matter and tachyon descriptions at that point.  The seed is the tachyon image of \(V3\) under \(J_B\).  It needs no gain, because
in the tachyon frame the major content sits in the partner port and the pair flux alone reaches \(S=20\ge16\).  The derivation
does not recover r65.  It shows that a calculation which fixes the frame factor at its Hopf value \(\gL=2\), as r20 proposed,
lands on the one point where matter and tachyon agree.

\subsection*{v4.9.0-r21 the bound in each frame}
\label{subsec:v490-r21-per-frame}

\begin{center}
\small
\begin{tabular}{@{}>{\raggedright\arraybackslash}p{0.33\textwidth}rrrrrl@{}}
\toprule
frame, point & \(|G|\) & \(\kappa\) & \(N\) & \(S\) & reserve & remark\\
\midrule
matter, Hopf & \(2\) & \(0.250\) & \(5\) & \(10\) & \(0.200\) & window\\
tachyon, Hopf & \(1\) & \(0.750\) & \(5\) & \(20\) & \(0.200\) & no window needed\\
matter, measured & \(1.864\) & \(0.232\) & \(5.31\) & \(10.63\) & \(0.1925\) & window\\
tachyon, measured & \(1\) & \(0.768\) & \(5.31\) & \(22.92\) & \(0.302\) & no window needed\\
photon, disk (\(b=0\)) & \(\infty\) & \(0.5\) & --- & --- & \(\to1\) & degenerate\\
photon, scale line (\(G=1\), \(N=5\)) & \(1\) & \(0.5\) & \(5\) & \(15\) & \(<0\) & blocked\\
\bottomrule
\end{tabular}
\end{center}

\paragraph{Theoremlet v4.9.0-r21.2.}
\begin{enumerate}[leftmargin=2em]
\item[(a)] Matter frame: at the measured offset the bound is \(0.1925\), \(0.0075\) below \(1/5\).  The frame-free \(1/5\) is slightly
too generous there.  The collar row of r05 (\(1/6\)) stays inside and the C5 stage (\(1/4\)) outside.
\item[(b)] Tachyon frame: the pair flux reaches the relay threshold without the clock.  \(W_{\rm clk}^{\rm loc}\) has no content
there.  The unassisted reserve is \(0.302\).
\item[(c)] Photon frame: the bound is degenerate.  In the state disk the gain \(1/|b|\) diverges and every drift is tolerated.
On the scale line of r17, with no frame factor, the window is closed.  This is the information singularity of r17.2 in the
clock layer: at the photon the quantity that the bound constrains is not defined.
\item[(d)] Away from the Hopf point the matter and tachyon reserves differ, \(0.19\) against \(0.30\).  A frame-free bound
exists only at \(x=1/3\).
\end{enumerate}

\paragraph{Proof sketch.}
Numerical evaluation of r21.1 and of the window condition at the stated values.  \(\square\)

\paragraph{Relation to r20.}
r20 listed several ways to carry the chart away from the Hopf point.  The readings with the chart held fixed gave \(0.10\) to
\(0.14\) and put the collar row outside.  They violate (D3).  Under (D1)--(D5) the balanced reading is selected, the matter bound
is \(0.1925\), and consequence 3 of r20 does not apply.

\subsection*{v4.9.0-r21 a rule for frame distinction}
\label{subsec:v490-r21-rule}

The computation suggests a rule for building the frame distinction into other places.  The tachyon frame is the \(J_B\) image
of the matter frame: \(x\to1/x\), \(p_\pm\) exchanged, \(|\gamma|\) unchanged and \(\gamma\) negative, the phase \(i\pi\) of the
superluminal rapidity of r20 (v4.8.0).  The photon frame is the fixed point \(p_+=p_-=\tfrac12\), where readout quantities
degenerate.  A statement that is valid in one frame needs a separate check in the other two.  It is frame-free only where the
descriptions agree.  For the clock bound that is the Hopf point alone.  This layer proposes the rule.  Applying it to the other
layers is left open.

\subsection*{v4.9.0-r21 guard status}
\label{subsec:v490-r21-guard}

\(W_{\rm clk}^{\rm loc}\) stays a local statement.  Its condition becomes frame-typed: \(\Delta_\gamma<\Delta_M^{\max}(\bsym^2)\) in the
matter frame, \(1/5\) at the Hopf point and \(0.1925\) at the measured offset; empty in the tachyon frame; undefined at the
photon.  The narrower matter-frame value claims less than \(1/5\).  No guard is spent.  The photon-frame degeneracy is no route
to \(W_\gamma^{\rm glob}\): it is the absence of the constrained quantity, not a global clock statement.

\subsection*{v4.9.0-r21 register entries}
\label{subsec:v490-r21-register}

\begin{center}
\small
\begin{tabular}{@{}>{\raggedright\arraybackslash}p{0.34\textwidth}>{\raggedright\arraybackslash}p{0.58\textwidth}@{}}
\toprule
statement & status after v4.9.0-r21\\
\midrule
\(\Delta_\gamma<1/5\) & derived under (D1)--(D5) as the unique frame-consistent bound, at the Hopf point (r21.1); r65 text still lost\\
clock bound per frame & matter \(0.1925\) (measured), tachyon no window (reserve \(0.302\)), photon degenerate (r21.2)\\
clock seed \((5,\tfrac34,1)\) & tachyon image of \(V3\) under \(J_B\)\\
r20 consequence 3 & superseded under (D3): collar row inside\\
frame distinction & rule proposed (tachyon = \(J_B\) image, photon = degenerate fixed point); application to other layers open\\
\bottomrule
\end{tabular}
\end{center}

\subsection*{v4.9.0-r21 ledger and guard update}
\label{subsec:v490-r21-ledger-guard}

\[
R_{\rm frame}^{v490\text{-}r21}=(0,0,1)
\]
for (bound derived under stated premises, values per frame computed, application of the frame rule elsewhere open).

The layer spends no guard:
\[
R_{\rm glob}^{(v490\text{-}r21,+)}=(0,1,1,1,1),
\qquad
\Gamma_{\rm guard}^{v490\text{-}r21}=(1,1,1,1,1),
\qquad
\beta_4=1/256\ \text{unspent}.
\]
It does not imply exact finite \(C5\) zero closure, exact all-readout terminality,
\(C6/T2\), a global photon-clock/RH theorem, or global all-readout closure.

